\documentclass[12pt,a4paper]{article}

\usepackage{fontawesome5}
\providecommand{\faPencilAlt}{\faIcon{pencil-alt}}

\usepackage{styles/couleurs}
\usepackage{styles/mise_en_page}
\usepackage{styles/boites}
\usepackage{styles/maths}
\usepackage{styles/figures}

\usepackage{enumitem}
\usepackage{array}
\usepackage{tabularx}
\usepackage{amsmath}
\usepackage{amssymb}
\usepackage{graphicx}
\usepackage{tikz}
\usepackage{xparse}

\usetikzlibrary{calc,arrows.meta}

\setlist[itemize]{leftmargin=1.4em,itemsep=0.15em,topsep=0.2em}
\setlist[enumerate]{leftmargin=*,itemsep=0.22em,topsep=0.25em}

\newcommand{\lignepointillee}{%
  \noindent\dotfill\par
  \vspace{0.45em}
}

\newcommand{\lignesreponse}[1]{%
  \begin{tcolorbox}[
    enhanced,
    colback=white,
    colframe=NeutreC!40,
    boxrule=0.5pt,
    arc=1.5mm,
    left=3mm,
    right=3mm,
    top=2mm,
    bottom=1mm,
    before skip=2pt,
    after skip=2pt
  ]
    \foreach \i in {1,...,#1}{\lignepointillee}
  \end{tcolorbox}
}

\newcommand{\imageactivite}[3]{%
  \IfFileExists{#3}{%
    \includegraphics[width=#1,height=#2,keepaspectratio]{#3}%
  }{%
    \fbox{%
      \parbox[c][#2][c]{#1}{%
        \centering\sffamily\small
        Image à intégrer\par
        \vspace{0.3em}
        \texttt{\detokenize{#3}}
      }%
    }%
  }%
}

\newcommand{\titreseq}[1]{%
\begin{center}
  {\sffamily\LARGE\bfseries\color{TitreC}
  Séquence 5 -- Statistiques et proportionnalité}\\[0.2em]
  {\sffamily\large\bfseries\color{SavoirC}#1}
\end{center}
}

\begin{document}

\titreseq{Activité 4 -- Quelles grandeurs faut-il utiliser ?}

\section*{Partie A -- L'exercice des fruits}

Un marchand vend des ananas à l'unité, à \(2\) € pièce.
Il vend des oranges au poids, à \(1\) € le kilogramme.

Bob achète deux ananas qui pèsent ensemble \(1{,}5\) kg,
ainsi que huit oranges qui pèsent ensemble \(2\) kg.

\begin{enumerate}
  \item Pour calculer le prix des ananas, quelle information est utile ?
  \item Quelle information sur les ananas est inutile ?
  \item Pour calculer le prix des oranges, quelle information est utile ?
  \item Quelle information sur les oranges est inutile ?
  \item Combien Bob paie-t-il ?
\end{enumerate}

\lignesreponse{7}

\section*{Partie B -- L'exercice du boucher}

Le rôti de bœuf coûte \(20\) € par kilogramme.

\begin{enumerate}
  \item Combien pèse un rôti qui coûte \(30\) € ?
  \item Combien pèse un rôti qui coûte \(18\) € ?
  \item Combien pèse un rôti qui coûte \(37\) € ?
\end{enumerate}

Présente les données dans un tableau de proportionnalité.

\lignesreponse{7}

\section*{Partie C -- Une maquette du Belem}

Le voilier \emph{Le Belem} mesure \(58\) m de long.
Une maquette mesure \(0{,}29\) m de long et \(0{,}17\) m de haut.

\begin{enumerate}
  \item Quel est le coefficient permettant de passer des dimensions de la maquette
        aux dimensions réelles ?
  \item Quelle est la hauteur réelle du Belem ?
  \item Le tirant d'eau réel est de \(3{,}60\) m.
        Quel est le tirant d'eau de la maquette ?
\end{enumerate}

\lignesreponse{7}

\aretenir{Avant de calculer, il faut identifier les grandeurs utiles,
leurs unités et le sens dans lequel on passe d'une grandeur à l'autre.}

\end{document}
